\SourceLeaf{18}
unde $r=\BareRoot{((a-x)^2+(b-y)^2+(c-z)^2)}$
\[
\begin{aligned}
\cos MX&=\frac{a-x}{r}\\
\cos MY&=\frac{b-y}{r}\\
\cos MZ&=\frac{c-z}{r}
\end{aligned}
\]
denique per theorema satis notum fit
\[
\cos MQ=\cos MX.\cos QX+\cos MY.\cos QY+\cos MZ.\cos QZ
\]
sive
\[
\cos MQ=\pm\frac{T(a-x)+U(b-y)+V(c-z)}{r\BareRoot{(TT+UU+VV)}}
\]

\Article{10}
\par
Iam ut integratio expressionum differentialium per totam superficiem absolvi possit, has expressiones ita transmutare oportet, ut duas tantummodo variabiles contineant. Hoc fieri quidem potest eliminando unam e variabilibus $x,y,z$ adiumento aequationis $W=0$: sed plerumque hoc modo formulae minus tractabiles prodeunt. Praestat itaque, duas novas indeterminatas $p,q$ introducere, ita ut tum $x$, tum $y$, tum $z$ tamquam functiones harum indeterminatarum considerare oporteat.

\begingroup\rightskip=0pt plus 1em
Simulac igitur ipsis $p,q$ valores determinati tribuuntur, etiam $x,y,z$ determinatae erunt, i. e. illis punctum determinatum in corporis superficie respondebit. Haec mutua correlatio clarius ob oculos ponetur, si planum indefinitum concipiamus, cuius singula puncta per coordinatas rectangulares $p,q$ exhibeantur. Cuivis itaque puncto plani respondebit punctum in superficie corporis et quidem unicum tantum, si res ita instructa est, ut $x,y,z$ sint functiones uniformes indeterminatarum $p,q$. Quodsi vice versa etiam per $x,y,z$ plene et absque ambiguitate determinantur $p$ et $q$, manifesto cuivis puncto superficiei corporis unicum tantum plani punctum respondebit, planumque in hoc casu undique in infinitum porrigi debet, quo integram corporis superficiem exhauriat. Alioquin autem plani partem tantummodo considerare oportebit, limitibus finitis vel infinitis descriptam, quae corporis superficiem quasi repraesentabit. Concipiatur planum per infinitas rectas tum lineae abscissarum parallelas tum ipsi normales
